\documentclass[a4paper, 12pt, openright]{ctexbook}

\usepackage{geometry}
\usepackage{hyperref}
\usepackage{fancyhdr}
\usepackage{titlesec}
\usepackage{enumitem}
\setlist{nosep} % 或者使用 nolistsep，彻底取消列表项之间及列表前后的额外间距
\usepackage{setspace}
\usepackage{xcolor}
\usepackage{tikz} % 引入 tikz 绘制标签
\usepackage{enotez} % 引入 enotez 处理尾注
% --- 引入图片处理宏包 ---
\usepackage{graphicx}
\graphicspath{{figs/}} % 建议设定统一的图片存放路径

% --- 1. 字体设置 (思源系列) ---
\setCJKmainfont{Source Han Serif CN}[ 
    UprightFont = Source Han Serif CN,
    BoldFont = Source Han Serif CN Bold,
    ItalicFont = Source Han Serif CN
]
\setCJKsansfont{Source Han Sans SC}[ 
    UprightFont = Source Han Sans SC,
    BoldFont = Source Han Sans SC Bold
]

% --- 2. 页面与行距 ---
\geometry{left=2.5cm, right=2.5cm, top=3cm, bottom=3cm}
\setlength{\headheight}{15pt} % 修复页眉高度溢出的警告
\linespread{1.5}

% --- 3. 元数据与链接设置 ---
\hypersetup{
    colorlinks=true,
    linkcolor=black, 
    filecolor=magenta,      
    urlcolor=blue,
    pdftitle={力量训练计划}, 
    pdfauthor={马克·瑞比托, 安迪·贝克},
    pdfsubject={健身},
    pdflang={zh-CN}
}

% % --- 4. 尾注设置 ---
% \setenotez{
%     list-name = {注释},
%     reset = true,
%     list-heading = \subsection*{#1} % 强制降级为 subsection，避免另起一页且缩小字号
% }

% --- 5. 出血标签 (统一灰色) ---
% 简化了原有逻辑，统一使用 gray 填充 
% \newcommand{\drawthumb}[1]{%
%   \ifnum\value{chapter}>0
%     \begin{tikzpicture}[remember picture, overlay]
%       \node[
%         fill=gray, % 统一替换为灰色
%         text=white,
%         font=\sffamily\bfseries\Large, 
%         minimum width=1.5cm, % [cite: 5]          
%         minimum height=2cm,  % [cite: 5]          
%         anchor=#1
%       ] at ([yshift={-2cm - (\value{chapter}-1)*2.5cm}]current page.#1) { \thechapter }; % [cite: 5]
%     \end{tikzpicture} % [cite: 6]
%   \fi
% }

% --- 6. 页眉页脚 ---
\pagestyle{fancy}
\fancyhf{}
\fancyhead[CE]{\leftmark} % 
\fancyhead[CO]{\rightmark} % 
\fancyfoot[C]{\thepage} % 
% \fancyhead[RO]{\drawthumb{north east}}
% \fancyhead[LE]{\drawthumb{north west}} 

% 覆盖 chapter 默认的 plain 页面样式，确保章首页也有标签，且不显示横线
\fancypagestyle{plain}{
  \fancyhf{}
  \fancyfoot[C]{\thepage}
  \renewcommand{\headrulewidth}{0pt}
  % \fancyhead[RO]{\drawthumb{north east}}
  % \fancyhead[LE]{\drawthumb{north west}}
}

% --- 7. 标题格式 ---
\ctexset{
    chapter = {
        format = \centering\Huge\bfseries,
        number = \chinese{chapter}, % [cite: 2]
        name = {第, 章}, % [cite: 2]
        beforeskip = 10pt, % [cite: 2]
        afterskip = 30pt % [cite: 2]
    },
    section = {
        format = \Large\bfseries,
        beforeskip = 20pt,
        afterskip = 15pt
    }
}

% --- 8. 内容层次自定义环境 ---
% % 文人简介环境：使用楷体，左右边距缩进，并强制首行缩进两字符
% \newenvironment{authorintro}
%   {\list{}{\rightmargin2\ccwd \leftmargin2\ccwd \itemindent=2\ccwd \listparindent=2\ccwd}%
%    \item\relax\kaishu}
%   {\endlist\vspace{1em}}

% % 文章前言 (题解) 环境：使用仿宋，整体字号略小，左右边距缩进，并强制首行缩进两字符
% \newenvironment{preface}
%   {\list{}{\rightmargin2\ccwd \leftmargin2\ccwd \itemindent=2\ccwd \listparindent=2\ccwd}%
%    \item\relax\small\fangsong}
%   {\endlist\vspace{1em}}

% 引用环境 (覆盖默认 quote)：缩进逻辑与前言保持一致
\renewenvironment{quote}
  {\list{}{\rightmargin2\ccwd \leftmargin2\ccwd \itemindent=2\ccwd \listparindent=2\ccwd}%
   \item\relax\kaishu\small\color{black}}
  {\endlist}

% --- 引入并配置标题宏包 ---
\usepackage{caption}
% 向 caption 宏包注册 \kaishu 命令
\DeclareCaptionFont{kaishu}{\kaishu}
\captionsetup[figure]{
    format=hang,               % 悬挂缩进：长标题换行后与首行文字对齐，不跑到标签下方
    labelsep=quad,             % 标签（如“图 1.1”）和标题文字之间用一个空格隔开
    font={small, stretch=1.2}, % 字号调小为 \small，局部行距缩小为 1.2（覆盖全局的 1.5，避免长标题太散）
    labelfont={bf, sf},  % 标签加粗，并使用你配置好的无衬线字体（思源黑体）
    textfont={kaishu},         % 标题正文使用楷体，与你现有的 quote 环境逻辑保持和谐
    justification=justified,   % 两端对齐，适合长段落阅读
    singlelinecheck=false      % 设为 false 强制所有长短标题都左对齐；若喜欢单行标题居中，可设为 true
}

\begin{document}

\frontmatter

\title{\Huge \textbf{力量训练计划}}
\author{\textbf{【美】马克·瑞比托\quad 安迪·贝克}}
\date{}

\maketitle

\tableofcontents

\mainmatter

\chapter{引言}

欢迎来到《力量训练计划》第3版。与之前有些许不同，新版本将从详细阐述什么不是力量训练着手来介绍力量训练的基本概念。有三个术语经常用于形容健身者的活动：“身体活动”（Physical Activity）、“锻炼”（Exercise）和“训练”（Training）。本书只讲“训练”，所以我们会首先定义前两个术语，“训练”一词则放在后面详加阐述。

\textbf{身体活动}是美国心脏病协会（The American Heart Association，简称AHA）建议你每周都要做的事情。“一切可以让你动起来消耗些热量的事情都可以称为身体活动”，该协会的网站上如是说。他们认为活动身体是延续生命所必需的。从本质上讲，除了坐着或躺着，其他形式都可以算是身体活动。本书不太关注这一点，因为即使是老年人，也能找到比在建议时段随便动动更有效的活动方式。

身体健康（Physical Fitness）是个相对概念。瑞比托和基尔格（Kilgore）在2006年的《运动生理学在线期刊》（\textit{Journal of Exercise Physiology Online}）［9（1）：1-10］中将其定义为：

\begin{quote}
    “与人类基因型的功能性性状表达一致，拥有能够为成功的事业努力、娱乐追求及家庭责任提供足够力量、耐力和灵活性的身体状态。”
\end{quote}

这个定义无疑是对前人量化该定义的改进：它既在整个生命周期内保持了框架的开放性，又是基于进化理论提出的，并且阐明了从遗传角度看待健康的必要性。根据这个定义，人类基因型的最优表达方式为：一个健康的人。这个定义从很多层面上讲都是令人信服的。

但是，我们并不打算止步于此，因为我们是运动员，因为我们热爱竞争——即使只能与自己竞争，我们不满足于仅仅是不坐着的“身体活动”，我们寻求优化“身体素质”的表达方式。我们力求提高身体素质，这在美国心脏病协会看来可能是过度的，更不是保持健康所必需的——毕竟他们想要的结果只是更多的人不会死于心脏病罢了。

说到这里，“锻炼”和“训练”这两个不同的概念更值得准确定义和体会。这两个术语常常被混用，显然这是不对的。“练习”（Workout）这个词常被用在锻炼和训练中，来指一个计划好的、能给身体施加一定刺激的事件（我们通常不会把汽车没油后推车的行为称为“练习”，尽管它也能产生同样的身体刺激）。锻炼和训练都要通过“练习”完成，但这两个概念却有着显著的差别。

\textbf{锻炼}是指为了当日、当时的效果所进行的身体活动。每次练习提供的身体刺激都是为了满足训练者当下的需求：消耗些热量、使身体暖和些、出出汗、喘喘气、让肱二头肌产生泵感、拉伸关节等等，基本与上下班打卡无异。锻炼就是为了练而练，不论是在练习期间还是练习结束后。锻炼的内容可能每次都是一样的，只要它能让你获得当时想要的感觉就行了。

但是对有着明确的运动表现目标的运动员来说，\textbf{训练}是必需的。在这个语境下，训练是为了实现长期目标而进行的身体活动，因此，训练是个过程，而不是组成这个过程的某些练习。并且，这个过程必须在一次次的练习之后、在某个时间点产生一个可见的结果，只有精心地计划这个过程才能取得这样的结果。训练的目标是运动表现的长期提高，这不仅需要时间，而且需要不达目的誓不罢休的主观意愿。

大多数人不是竞技运动员，也不会把自己看作是竞技运动员，除了减肥、塑形外也没有什么明确的目标，这与不考虑基因和表型因素时保持身体健康的要求很接近。所以，大多数人非常满足于锻炼。健身行业不仅深谙此理，而且乐意迎合。因此，缺少杠铃训练的系统性加重组件和平衡性、基于训练机的练习成了主要的锻炼方式。P90X、CrossFit以及各种DVD上声称能“迷惑”你肌肉的乱七八糟的计划也是如此。现代商业健身房基本上是为锻炼而生的，因为训练会让它们无利可图。标准的商业健身房会在55\%的地面空间放置那些人们可以看着电视完成各种重复动作的有氧器械，剩余45\%的空间则放置那些主要为方便工作人员而设计的固定器械，因为它们用起来简单、教起来容易、清理起来不费事。

很多商业健身房除了哑铃没有任何其他的自由重量，它们也无意教你使用自由重量进行锻炼。这些健身房本质上是商业组织，不是锻炼的地方，如果会员办了卡来了两三次就不再来了，对他们来说也无所谓。他们的商业模式建立在大量雇佣体育类专业（运动生理学、生物力学或者别的什么专业）的大学生作为廉价劳动力的基础上，这些学生只是一群没有任何杠铃运动的执教经验，也没有读过这本书的孩子罢了。这种商业模式依赖于运动人群的快速流转，尤其是在高峰期，如同健身房规划的那样，主要都是些使用有氧器械的人——也就是那些需要依靠这些俱乐部锻炼的人。俱乐部则期望会员会在固定器械上玩上20分钟，在单车或者跑步机上看着电视活动半小时，然后洗个淋浴拍拍屁股走人。员工的职责就是让这个过程高效执行。在这种情况下，锻炼很容易完成，训练则几乎不可能实现。

目前，健身行业正在经历剧烈的变革，希望是在向好的方向改变。55\%的有氧器械和45\%的固定器械的标准行业模式正在让位于“功能性训练”设施。这种训练强调强度更高的杠铃动作、全身性的体操动作，并与那些更加费力且真正能够给身体施加足够刺激以引发适应状态的运动结合起来。这与本身发展迅速却禁止任何高强度训练、禁止发出噪声的“星球健身”（Planet Fitness）形成了鲜明对比（译者注：星球健身是美国一家连锁健身俱乐部，因为担心会员被吓到而禁止大块头健身爱好者入内，禁止大重量训练，缺少自由重量，只有固定器械和大量有氧器械，其中最具特色的一点是当有人使用大重量训练喊出声或使用器械发出太大声响时会响起警报）。但很多功能性训练设施的问题在于工作人员准备不足。这些工作人员一般是一些热情饱满却缺乏执教经验的年轻人，因此他们无法保证会员以正确的方式进行训练并避免受伤。而且，即使是CrossFit风格中动作的难度很高，我们仍然不能称之为训练。

训练要花费时间，要有人指导，要为了目标不懈努力。这就需要一个计划——一个由熟悉这个过程和了解完成这个过程需要付出多少努力的人来制订的计划，以及一种意愿——愿意接受每种动作只有处在一系列能够达成最终目标的事件中才有价值的事实。

但这并不意味着每次的训练都枯燥无味，更不意味着训练中的渐次进步无法赋予运动员渴望的成就感。我的意思是，每个运动员都要明白，每次训练都是为了完成最终画幅的一块小小拼图，而这幅最终拼图的美丽，是那些只知道锻炼而缺少整体规划的人永远无法想象的，他们甚至对此一无所知。

\section{力量训练}

这本书的主题是训练而不是锻炼，并专注于力量训练。不论是什么项目，运动员想要取得成功都必须训练，我们的力量训练与长跑训练的理念是相同的：从运动员现有的水平起步，为了将要达成的目标做计划。只不过我们主要利用无氧阻力训练，而长跑运动员主要利用有氧耐力训练。但训练对我们有着同样的意义——我们走进举重室不是为了拿起杠铃玩耍，他们走上跑道也不会乱蹦乱跳直到感觉无聊。训练意味着计划，而制订计划需要了解我们想要改变运动员哪方面的身体能力。

力量训练是一个提高运动员的肌肉力量以对抗外部阻力的过程。力量训练遵循正确的逻辑进程，从运动员当前的力量水平起步，向着增强力量的目标前进。这样的进步需要两个条件。

第一，如果你想为提高运动员的力量水平制订计划，对其当前的力量水平进行正确的评估是必要的。这项评估可以也应当在教授运动员完成计划所需的举重动作时进行，因为这些动作必须正确地传授和完成。学习举重的过程需要逐渐增加重量，因为我们想要变得更强壮就必须能够举起更大的重量。在启动训练计划时，进行专门的评估不仅浪费时间，更重要的是，这样做忽视了评估过程本身也是形成身体适应状态的刺激因素。如果测试强度足以精确地评估运动员当前的身体能力，那么由评估过程产生的对刺激的适应状态已经改变了评估对象的身体能力，导致提供的数据不准确，失去了为测试后的任意时间点提供参考的价值。在第一次举重教学课上边教学边评估运动员的身体能力则要高效得多。对所有的初级训练者来说，训练计划要以此为始，第二次训练则以第一次训练结束后的水平为起点，这样教学和评估的目的就能同时达到了。

第二，训练计划的制订必须以最高效地帮助运动员增长力量为目标。我们即将会看到，这需要用到一些基于刺激—恢复—适应模型的基本训练法则，同时也需要对运动员当前的身体适应潜力进行正确评估。

\textbf{刺激}指的是任何可以改变机体生理状态的事件。刺激可以是一次高强度的运动，可以是晒太阳，可以是被熊暴揍了一顿，或者3个月卧床不起。刺激会打破\textbf{体内平衡}（机体内的正常生理环境）。之后，为了更好地生存，机体会从刺激中\textbf{恢复}到比施加刺激之前的状态更强一点的状态（如果能够恢复的话——被太阳晒黑了会很容易恢复，但熊的袭击就是个问题了），以应对同样的刺激再次发生的情况。机体对于刺激的这种\textbf{适应}是其保证自身在多变的环境下存活的方式。而且，这种对刺激的适应过程是生命的重要特征之一。

本书中的刺激是通过认真地使用杠铃训练产生的，并借此使身体产生促进肌肉力量增长的适应状态。与机体所承受的任何其他类型的重复性刺激一样，身体对先前刺激的适应会积累，进而从根本上改变机体。很显然，与刚出生时相比，你的生理状态已经大不一样了。除了正常的生长发育，这种变化也源自你在这些年里所经受的各种刺激。

在我们的训练中，我们可以持续施加在身体上的刺激类型与你之前经受过的刺激有关，因为你目前的适应状态构成了身体对刺激的最终适应潜力的一部分。不论是紧急状况下的急性刺激，还是在一段时间内产生的慢性刺激，每个人适应刺激的能力都存在上限。这个上限由基因天赋和运动员所处的物理环境决定，这两点共同控制着每个人的运动表现的潜力。其实，人类所有方面的潜力上限都被这样的进程所控制，这也解释了为什么每个领域的顶尖人才都很少见的状况。图\ref{fig:1-1 力量表现提高与训练复杂度相对于时间的一般关系}总结了这些概念。

\begin{figure}[htbp] % htbp 控制图片排版位置优先顺序
    \centering
    \includegraphics[width=0.8\textwidth]{image_000.jpg}
    \caption{力量表现提高与训练复杂度相对于时间的一般关系。注意，对训练的适应速率会随着训练生涯的延伸逐渐放缓}
    \label{fig:1-1 力量表现提高与训练复杂度相对于时间的一般关系}
\end{figure}

每个人接近其上限的程度取决于这个人还有多少潜力可供挖掘。比如，从发掘身体潜力的角度看，一个未经训练的17岁少年和一个38岁的竞技型举重运动员可以说是两个极端。这个少年的力量潜力几乎没有得到开发，而竞技型举重运动员已经致力于提高力量20年了，他已经十分强壮了；这个少年的所有潜力几乎都可以得到挖掘，而举重运动员已经把所有的潜力最大限度地转化为了自身的能力；这个少年能够很快、很容易地变强壮，而这位举重运动员用几个月的时间执行一项十分复杂的训练计划也只能稍稍提高一点力量，因为他已经非常强壮了。如果你还不够强壮，那么变强壮会相对容易一些。这个少年每次训练之后取得的进步可能比这个运动员半年时间取得的进步都要大。你或许会觉得这很悲哀，或许觉得很奇妙，看你思考问题的角度了。

与无数的自然现象和社会现象一样，人类的运动表现符合收益递减原则（Principle of Diminishing Returns）。试着加速物质到光速、学习钢琴、打造一辆速度更快的汽车等，这些事情都是开始时比较容易，逐渐变得越来越困难，你需要投入越来越多的能量、时间或金钱才能接近极限，但不可能达到真正的极限。人类的运动表现不也是如此吗？一开始很容易，接近最终目标则很难——从来没有人的深蹲力量能够在1年内增长200磅（90.7千克），否则每次比赛时世界纪录都会轻易被打破了。

这个道理显而易见，但大多数训练计划都忽视了这一点。我们学到的一直是让初级训练者测试各种项目的单次最大重量（1 Repetition Maximum，简称1RM），但初级训练者不懂得要如何完成动作，也不懂得如何才能做得正确，因此他们不能做得足够好，这样最终测出来的1RM并不具有实际意义。然后，基于这些糟糕的数据，健身专家会给初级训练者安排一个实际上更适合高级运动员的计划：初级训练者需要在大部分的时间里使用次大重量，并按照预先的设定加重，有时甚至会每月增加一次重量。这样的计划不能正确地匹配初级训练者在训练初期可以快速适应的能力，正是因为之前没有经历过这样的快速适应，所以他们在训练初期才有能力做到这一点。

更糟糕的是，随着时间的推移，初级训练者可能因此无法取得任何有意义的进步。按照上述方式训练，教练通常会建议，当你使用能够完成4～5组、每组8～12次重复的重量感觉更容易时，你就可以增加一点重量了。这种方式缺少任何能驱动进步的尝试。但如果你碰巧这样做取得了进步，只要不受伤，也不是不可以。

这种“传统智慧”在训练计划的制订中非常常见，甚至很多认证机构把其中的一些版本看作正确的、合适的。这些机构包括美国运动医学会（America College of Sports Medicine，简称ACSM）、美国国家体能协会（National Strength and Conditioning Association，简称NSCA）、理念健身（IDEA）、美国运动协会（The America Council on Exercise，简称ACE）、美国田径与健身协会（Athletic and Fitness Association of America，简称AFAA）、美国运动和体育训练协会（National Exercise and Sports Trainers Association，简称NESTA）、美国体育与健身协会（American Sports and Fitness Association，简称ASFA，除非你通过考试，否则不收费！）、基督教青年会（Young Men's Christian Association，简称YMCA）、库珀诊所（The Cooper Clinic）。因为这些计划是有依据的，也就是说有很多经过同行审查的运动科学文献支持你这么做。

但锻炼不是训练，我们制订训练计划的目的与他们不同，我们想要的是通过训练让人们变得更强壮，而他们的主要目的是培养私人教练和运动课程讲师。

\section{理论方法}

在我们的训练方法中，每个训练者的训练计划都必须根据他们当前在图1-1的曲线中所处的位置来制订。根据训练者从训练引起的体内平衡破坏中恢复的快慢不同，本书采用“初级”“中级”和“高级”三个术语来描述训练者的等级。这三个术语并不是指训练者的力量水平或者绝对运动能力。这些术语对于不同的运动项目有着不同的应用意义，但在本书中特指图1-1模型中对训练者的划分。

由于初级训练者之前从未按照有规律地逐渐增长力量的计划训练过，所以他们举起的重量相对于其最终的力量和爆发力潜力来说要小得多。这里的初级训练者可能有过多年健身房健身的经历，每周去健身房且风雨无阻，但他们从来没有真正训练过。实际上，初级训练者在一次训练后只需要48～72小时就能完全恢复，他可以周一使用所谓的“大重量”，然后周三再次进行“大重量”训练。这些训练者远未达到其身体潜力的上限，他们的力量和神经效率尚未得到充分锻炼，因此他们无法使用足以影响其恢复速率的大重量。实际上，他们使用的“大重量”并不真的很大。在力量和爆发力增强的同时，你的恢复能力也在提高。恢复能力和其他的身体参数一样是可以训练的，并且还是训练过程中极为重要的环节。但在现实的训练中，往往会出现不当的或者过量的训练刺激超越训练者当前恢复能力的情况。你要记住，想要进步，必须先恢复。

在这里，我们把\textbf{初级训练者}简单定义为：在下一次训练前能够从上一次训练所施加的刺激中恢复过来并足以产生适应状态的训练者。这就意味着在整个初级阶段，训练者每次都能够增加正式组的重量，从而在较短的时间内实现力量的快速增长。如果训练由一个懂得初级阶段的过程及其潜在效能的人好好安排的话，初级阶段会成为运动员训练生涯中力量和能力进步最快的阶段。

当力量的增长出现平台期时，初级阶段就结束了。平台期往往出现在训练的第三个月到第九个月，具体时间取决于每个人的基因天赋以及对影响恢复的环境因素的控制程度。初级阶段的训练计划基本上就是我们在《力量训练基础》中专门为自由重量训练设计的线性模型。

有一个重点需要好好理解——初级训练者适应于不活跃的状态（与重量训练相关），所以即使使用那些不是专门为了增强力量而设计的训练计划来完成基础杠铃动作也能获得力量的增长。比如，每次练习几组20次重复的杠铃动作同样能促进初级训练者的1RM的增长。一个之前整日久坐不动的初级训练者甚至骑单车都能增加其深蹲的1RM。但对中级训练者和高级训练者来说，他们在力量、爆发力和增加肌肉量方面的提升都与特定训练计划的合理运用密切相关。

初级训练者每次训练都会完成两件事：用全新的、更大的重量来“测试”自身的力量水平，并且这个测试重量会使身体在下次训练到来时变得更强壮。按照预定的组数和重复次数训练，每次比上次增加10磅（4.5千克）重量，这既证明了上次的训练对促进力量增长是成功的，又促使身体适应这种刺激，在下次训练到来时变得更强壮。对绝大多数训练者来说，只要训练组织得当，他们会在初级阶段取得整个力量训练生涯中最快速、最有成效的力量进步。

当训练者取得快速且轻松的进步开始变得困难、每次训练取得进步变得越来越不容易时，初级阶段就接近尾声了。更小的重量增幅将被使用并会很快用尽，不论如何想方设法地恢复，进度都会停滞下来。

\textbf{中级训练者}则要解决一系列不同的问题。中级训练者开始使用更接近自身身体潜力的重量，他们的恢复能力受到外界刺激的影响也不同。恢复需要更长的时间——而这段时间内还要进行多次训练。从实际出发，使用“周计划”安排训练最为高效。中级训练者已经发展出了承受恢复周期更长的刺激的能力，同时，这个阶段打破身体平衡所需要的刺激已经开始超出先前阶段只须48～72小时恢复的能力。想要同时兼顾足够的刺激和充分的恢复，训练负荷必须在一个更长的时间周期内变化，以周为单位安排训练计划是种便利的做法。开始时，恢复所需要的实际时间可能少于1周，比如5天，到了中级训练阶段的末期，恢复时间可能会增加到8～9天。要点是该如何分配增加的训练量，从而在施加足够刺激的同时又能让恢复不那么困难。在这个阶段训练成功的关键是，要平衡这两个重要却相互影响的要素——不断增长的刺激需求和相应增加的必要恢复时间。在一个单独的训练周期内，简单地按周规划一次或多次大重量训练有助于训练者随后的恢复，并且这样的规划与日历的日程安排相匹配。

相对初级训练者来说，中级训练者能从多样的训练动作中受益更多。这个阶段的运动员会通过学习新的动作模式来发展自己的技术，这样做的同时，他们也提高了自己获取新技术的能力。在这个阶段，训练者真正成了运动员，可以选择一项运动，做出一些会影响整个竞技生涯的决定。如果他们广泛地接触到了各种各样的训练方式和竞技运动的话，他们能够更有效地做出决定。

随着一系列难度逐渐增加的周计划的完成，你会遇到运动表现的平台期，这标志着中级训练阶段走到了终点。取决于个体对全年性渐进训练的耐受性和执行度，中级阶段短则2年，长的话可以持续4年，甚至更久。其实有75\%甚至更多的训练者很可能永远用不到比中级计划更复杂的训练计划——请记住，训练者处于什么阶段并不取决于能够举起的重量或者训练的年限。基本上，所有的非杠铃项目的竞技运动员的力量训练可以一直采用中级训练计划。这些运动员的训练不会局限在举重室里，他们会更专注于专项竞技项目的训练。这会有效地延长中级训练阶段持续的时间，在这个过程中，即使是非常成功的运动员也可能永远无法榨干中级力量训练计划带给自己的好处。

杠铃项目的\textbf{高级训练者}，其训练强度已经接近于其最终的力量潜力了。这个级别的训练所需要付出的时间和精力不是一般的训练者可以达到的，这类人是训练人群中的少数，几乎完全由竞技型的力量举和奥林匹克举重运动员构成。由于运动员的恢复能力是可以通过训练提高的，所以高级举重运动员所承受的训练负荷相当高。不过，由于在进阶到高级阶段的过程中身体已经产生了的适应状态，所以高级运动员产生足够的身体适应所需要的训练负荷也会相当高。这个级别的训练量和训练强度非常大，所以需要比使用中级训练负荷时更长的恢复时间。高级阶段的训练需要使用更复杂、更多变的方法，并在更长的时间周期内安排训练负荷和与恢复相关的参数。把二者结合起来，兼顾能够成功取得进步的训练负荷和身体恢复的周期往往需要一个月甚至几个月。比如，我们可以安排1周的大重量训练来打破体内平衡，而这周的训练后可能需要3周甚至更长的时间使用较小的重量训练才能使身体完全恢复并产生适应状态。这个阶段的进步曲线非常平缓（见图\ref{fig:1-1 力量表现提高与训练复杂度相对于时间的一般关系}），以一种非常缓慢的速率逐渐接近身体的最终潜力，此时为了取得一点点的进步都需要付出极大的努力。同样因为这个原因，高级训练者使用的训练动作的数量通常也要少于中级训练者。因为高级训练者已经适应并专精于他们特定的竞技项目，不再需要新的动作模式和刺激类型了。

控制训练参数的复杂性适用于高级训练者。大多数训练者永远不会达到需要安排周期化的高级训练计划的级别，因为他们往往会在达到高级训练阶段之前主动结束自己的竞技生涯。

\textbf{精英级运动员}指的是那些已经达到所在运动项目中“精英”标准的运动员。按照这个定义，参加国家级或国际级比赛的中级运动员也可以获得“精英”称号。当偶尔出现一些富有进取心且极有天赋的运动员时就会出现这种情况。我们都见过这样的基因怪物，他们看起来没有像其他人那样付出那么多努力，却能在一项运动中迅速崛起。精英级运动员通常是高级运动员的一个特殊子集。精英级运动员是极少数天赋极佳、恰恰又充满争取成功的动力，并且无惧体力成本和社交成本而奋勇前进的人。他们由于自身的成功而一直从事某种运动，尽管在训练中得到了巨大的回报，但他们仍然致力于这个级别的训练。

之前的训练把运动员带到了十分接近自己身体极限的水平，想要取得额外的进步就需要更加复杂的训练计划，才能把尚未兑现的些许潜力一点点挖掘出来。这些运动员必须使用非常复杂的训练计划——刺激是高度可变的，但动作选择可能很简单——逼迫适应状态已经非常强的自己更靠近终极的运动表现。在这个级别，训练计划可能会持续几个月、1年，甚至是奥运会的4年周期。这个级别的运动员使用的任何训练方法都是高度个性化的，也超出了本书的讨论范围。不论训练经历如何，能达到这个级别的运动员在所有训练者中的占比远低于0.1%。

与初级训练者和中级训练者不同，高级训练者需要大量的高强度训练才能打破体内平衡，强制身体产生新的适应状态。这就意味着，进步所需的刺激会慢慢接近身体可以承受并能从中恢复的最大负荷。此外，对高级运动员来说，进步的窗口已经非常小了。一个能通过完成10组深蹲训练取得进步的高级运动员，如果完成9组可能不会取得任何进步，而做11组又可能导致“过度训练”。就是这样！

但是如果训练负荷没有增加，体内平衡就不会被打破，运动员也就无法取得运动表现和恢复能力的提高。如图\ref{fig:1-1 力量表现提高与训练复杂度相对于时间的一般关系}所示，增加训练负荷的方式取决于训练水平。初级训练者对训练的适应与中级、高级训练者大不相同，根据每个训练阶段的不同生理指标来制订各个阶段的训练计划极为重要。如果你不想在举重室里浪费时间的话，就不要把一个高级举重者的训练计划扔给一个新手。

\section{有疑问？}

如果真相真如我说的那样显而易见，为什么那些认证机构不认可这样的模式并相应地改变他们的教条呢？会不会是那些掌控着训练“传统智慧”的学术组织还没有研究过这个模式，因此没有将其发表在由守护“传统智慧”的同行进行评审的期刊上呢？

我们以更好的方式重新梳理下这个问题：为什么那么多每年培养大量体育教育学士，以及众多硕士和博士研究生的大学，却无法好好研究一个执行时间必须跨越多年，研究对象为学术部门无法接触到的、高度积极的、不愿意为此改变个人训练以比较不同方法优劣的竞技运动员群体的课题呢？

答案很显然：他们做不到。大学的研究机制限制太多，无法设计出一个能够有效观察、对比不同训练方法对运动员影响的研究方案。学术部门能接触到的研究对象就是他们的学生，而这些学生基本上都是初级训练者水平，新手效应会响应任何的训练，也就是任何训练计划或多或少都有效。学术部门也能用老年人做研究对象，因为老年人有时间来配合研究。但学术部门很难找到竞技型的运动员，让他们改变自己的训练方式去迎合这些研究——由一些并不了解他们的运动专项和训练需求的人设计的研究方案。一个学期3个月，论文发表则是按年度的。硕士研究生一般会学习2～3年。他们必须发表论文来毕业，或者按照学位评审委员会主席的要求来写论文，而主席的工作是在管理层面前展现工作成效。负责研究设计和研究方法制定的人必须亲身体验过正确制订的训练计划才能设计出正确的问题。说出来你可能不会相信，在体育教育部门，满足这些条件的人凤毛麟角。体育教育部门的人要么忙于毕业，要么忙于争取终身职位，要么绞尽脑汁少教几堂课，要么想尽可能多发论文，要么已经到了退休年龄。这个评价听起来很刻薄，但并非指责这些人很邪恶。现实是，绝大多数的体育教育计划既不涉及训练概念，也不收集具体的训练数据，所以无法让研究者处理这些事情。

结果导致关于训练的文献存在极大的空白，无数由同行审阅的有关锻炼的文章反而充斥其中。因为锻炼的本质决定了体育教育部门更容易去研究它，训练则不行。只要目前的研究体系不改变，这个状况就不会得到改善。

以讲述锻炼为主的、经同行审阅通过的文献，形成了所谓的“有据可循的实践方式”——这个时髦的术语只能代表那些基于同行审阅的运动科学文献设计出来的锻炼方案，实际上根本不可能对真正的运动员训练提供什么帮助。通过研究少量普通群众的锻炼数据试图推论出适合运动员训练的结论，这种做法从未见效，也不可能有效，即使所有由同行评审的期刊都认同它，也无法改变其无用的事实。

经验主义是一种认为知识来自直接的感官体验——即经验证据——的认识论的观点。有些人会把经验证据看作在正式的研究环境中通过控制实验获得的数据。他们往往是制造这种数据的人，他们会把实验数据的缺失看作知识的缺乏。形成鲜明对比的是，理性主义（Rationalism）是一种竞争认识论，将理性和逻辑分析看作充分检验知识和真理的手段。对一个擅长进行理性分析的人来说，缺乏实验数据不是一个无法逾越的障碍，因为他们可以从普遍性原则推断出事情的细节。

来自经验丰富之人的观察——这里就是指那些有数十年执教经验、执教过数以千计运动员的教练的观察——经常被运动科学领域的发行机构、学术机构视为逸事、传闻，甚至是无稽之谈。这是对经验的误解，经验绝对应该包括富有经验的、了解情况的教练的直接观察。来自实验研究的证据只是经验证据的一种，它们同样来自于观察，与教练在实践中观察运动员采用的方式是完全相同的。因此，教练通过观察取得的证据同样具有价值，尤其是考虑到实验研究的数据好坏也受到数据获取方法优劣的影响。

运动科学研究自身存在问题。研究对象的样本通常非常小，每个实验组经常不足20人。在这些人中，运动员就更少了，一般都是些未经训练的大学生。对他们来说，任何刺激都能促使其产生适应状态。这种糟糕的研究方法无法对比两种不同训练方法的效果，并把与训练相关的所有问题完全排除在外了。研究方法本身的问题往往非常大（比如用史密斯机做深蹲研究），完全忽略了所研究的动作模式中的量化指标（确切来讲，什么是深蹲，深蹲深度要多大，髋角是多少，它对肌肉动员会产生什么影响，要如何测量，等等），抑或是研究人员无法用标准化的语言与研究对象交流（诸如“这次要非常非常地用力！”这样的表述）。有时，因为研究对象都是学生，他们只能配合一学期的时间，导致研究周期过短，无法就研究问题获得任何有意义的数据。最重要的是，如果主导研究的人毫无经验，无法判断课题本身是否存在问题（诸如“是躺在水平长凳上还是平衡的瑞士球上能卧推的重量更大”），同时审查人员也毫无经验，搞不清问题是否具有研究价值，最终，经同行审阅的“有据可循”的研究会成为新的参考文献，使这种糟糕的状况进一步恶化。

对长期的杠铃运动生涯来说，完成多组、每组5次重复的训练方式是最有利于发展力量的组数和次数的组合——这是基于观察证据得出的结论，可以使你免受先前的观念以及双盲对照研究的经验偏见的影响。二者都有自身的局限性，但也都有可取之处。或许根本不存在所谓的理论中性观察，但在缺少其他数据的情况下，来自了解情况的教练的观察就是我们拥有的最好的数据，根据这些数据得出的结论要远比那些糟糕的运动研究得出的推论可靠。在缺乏实验数据的情况下（就像减重和大腿肥大研究以外的训练方法面临的悲惨困境那样），经验主义和理性主义结合可以产生最好的结果。

由于缺乏杠铃训练长期效应的相关研究数据，我们不得不依赖成百上千的、在错误中摸索方向、总结经验的教练和运动员的观察结果。因此，每个能制订出有效的杠铃训练计划的人都是理性主义者。如果你想制订一份有逻辑、有效且高效的（这就是理性）杠铃训练计划，就必须具备坚实的生理学、化学和物理学知识，因为我们已经知道所谓的“运动科学”缺乏制订训练计划所必需的严格性和包容性。一个准备充分的教练要么有一个“过硬的”理学学位，要么有生物学、解剖学、生理学、物理学、化学或者心理学的背景。这些科目的教科书应当构成教练图书馆的基础，加上多年使用杠铃训练的实践经验、成千上万小时的执教经验，共同构成了其作为杠铃训练教练的能力。

\end{document}